% from ega4-2-fr.tex
\oldpage[IV-2]{25}
et si $\{u'\}$ est une base de $M$, il existe $v\in C$ tel que
$u'=u\circ v$, ce qui prouve que $u$ est surjectif~; donc \emph{a}$'$)
entraîne \emph{b}$'$), et comme \emph{a)} entraîne évidemment
\emph{a}$'$), cela achève de prouver le corollaire.

\begin{proposition}[2.5.8]
\label{IV.2.5.8-fr}
Soient $A$ un anneau commutatif semi-local, $B$ une $A$-algèbre (non
nécessairement commutative), $V$ et $W$ deux $B$-modules. Soit $A'$ une
$A$-algèbre commutative qui est un $A$-module fidèlement plat~; on pose
\[
B'=B\otimes_A A',\qquad V'=V\otimes_A A',\qquad W'=W\otimes_A A',
\]
de sorte que $B'$ est une $A'$-algèbre, $V'$ et $W'$ des $B'$-modules. On
suppose en outre vérifiée l'une des conditions suivantes~:
\begin{enumerate}
\item[(i)] $A$ et $A'$ sont noethériens, $V$ et $W$ sont des $A$-modules de
type fini.
\item[(ii)] $B$ est un $A$-module de type fini, $V$ est un $B$-module
projectif de type fini et un $A$-module de présentation finie.
\end{enumerate}
Alors, si $V'$ et $W'$ sont des $B'$-modules isomorphes, $V$ et $W$ sont des
$B$-modules isomorphes.
\end{proposition}

Notons que dans le cas (ii), $W'$ étant $A'$-isomorphe à $V'$, est un
$A'$-module de type fini, d'où résulte que $W$ est un $A$-module de type fini
(Bourbaki, \emph{Alg. comm.}, chap.~I\textsuperscript{er}, \S~3,
n\textsuperscript{o}~6, prop.~11)~; donc dans tous les cas $V$ et $W$ sont
des $A$-modules de type fini. En outre~:

\begin{env}[2.5.8.1]
\label{IV.2.5.8.1-fr}
Sous l'une des hypothèses (i), (ii), $\operatorname{Hom}_B(V,W)$ est un
$A$-module de type fini.
\end{env}

C'est évident dans le cas (i), car $\operatorname{Hom}_A(V,W)$ est alors un
$A$-module de type fini, et $\operatorname{Hom}_B(V,W)$ est un sous-$A$-module
de $\operatorname{Hom}_A(V,W)$. Dans le cas (ii), $V$ est facteur direct d'un
$B$-module libre $B_s$, donc $\operatorname{Hom}_B(V,W)$ est facteur direct de
$\operatorname{Hom}_B(B_s,W)=W^s$, et comme $W$ est un $A$-module de type
fini, il en est de même de $\operatorname{Hom}_B(V,W)$. Posons
\[
C=\operatorname{End}_B(V),\qquad M=\operatorname{Hom}_B(V,W)
\]
qui sont des $A$-modules de type fini dans les cas (i) et (ii). On sait que
sous l'une des conditions (i), (ii), l'homomorphisme canonique

\begin{env}[2.5.8.2]
\label{IV.2.5.8.2-fr}
\[
\operatorname{Hom}_A(V,W)\otimes_A A'\longrightarrow
\operatorname{Hom}_{A'}(V',W')
\]
\end{env}

est bijectif (Bourbaki, \emph{Alg. comm.}, chap.~II, \S~2,
n\textsuperscript{o}~10, prop.~11). Comme $A'$ est un $A$-module plat,
$\operatorname{Hom}_B(V,W)\otimes_A A'$ s'identifie canoniquement à un
sous-$A'$-module de $\operatorname{Hom}_A(V,W)\otimes_A A'$. L'image de ce
sous-module par l'homomorphisme (2.5.8.2) est contenue dans
$\operatorname{Hom}_{B'}(V',W')$, car si $u\in\operatorname{Hom}_B(V,W)$ et
$a'\in A'$, l'image de $u\otimes a'$ par (2.5.8.2) est l'homomorphisme
$u':V'\to W'$ tel que $u'(x\otimes1)=u(x)\otimes a'$~; pour tout $b\in B$,
on a donc
\[
u'((b\otimes1)(x\otimes1))=u'(bx\otimes1)=u(bx)\otimes a'
=bu(x)\otimes a'=(b\otimes1)(u(x)\otimes a'),
\]
d'où notre assertion. Cela étant~:

\begin{env}[2.5.8.3]
\label{IV.2.5.8.3-fr}
Sous l'une des hypothèses (i), (ii), l'homomorphisme
\end{env}

\begin{env}[2.5.8.4]
\label{IV.2.5.8.4-fr}
\[
\operatorname{Hom}_B(V,W)\otimes_A A'\longrightarrow
\operatorname{Hom}_{B'}(V',W')
\]
est bijectif.
\end{env}

Pour tout $b\in B$, notons $h(b)$ (resp. $h'(b)$) l'homothétie $x\mapsto bx$
dans $V$ (resp. $W$), qui est un $A$-endomorphisme. Soit
$(b_\alpha)_{\alpha\in I}$ un système de générateurs de la $A$-algèbre $B$~;
l'application
\[
u\longmapsto(h'(b_\alpha)\circ u-u\circ h(b_\alpha))_{\alpha\in I}
\]
de $\operatorname{Hom}_A(V,W)$ dans $(\operatorname{Hom}_A(V,W))^I$ est
$A$-linéaire, et par définition, son noyau n'est autre que
$\operatorname{Hom}_B(V,W)$~; autrement dit, on a une suite exacte
\[
0\longrightarrow\operatorname{Hom}_B(V,W)
\longrightarrow\operatorname{Hom}_A(V,W)
\longrightarrow(\operatorname{Hom}_A(V,W))^I.
\]
Le même raisonnement s'applique en remplaçant $A$, $B$, $V$, $W$ par $A'$,
$B'$, $V'$, $W'$~; de plus, on a un diagramme

\begin{env}[2.5.8.5]
\label{IV.2.5.8.5-fr}
\[
\xymatrix{
0 \ar[r] & \operatorname{Hom}_B(V,W)\otimes_A A' \ar[r] \ar[d]^{r} &
\operatorname{Hom}_A(V,W)\otimes_A A' \ar[r] \ar[d]^{s} &
(\operatorname{Hom}_A(V,W))^I\otimes_A A' \ar[d]^{t}\\
0 \ar[r] & \operatorname{Hom}_{B'}(V',W') \ar[r] &
\operatorname{Hom}_{A'}(V',W') \ar[r] &
(\operatorname{Hom}_{A'}(V',W'))^I
}
\]
\end{env}

où $r$ est l'homomorphisme (2.5.8.4), $s$ l'homomorphisme (2.5.8.2) et
$t$ l'homomorphisme composé
\[
(\operatorname{Hom}_A(V,W))^I\otimes_A A'
\xrightarrow{w}(\operatorname{Hom}_A(V,W)\otimes_A A')^I
\xrightarrow{s^I}(\operatorname{Hom}_{A'}(V',W'))^I
\]
$w$ étant l'homomorphisme canonique (Bourbaki, \emph{Alg.}, chap.~II,
3\textsuperscript{e} éd., \S~3, n\textsuperscript{o}~7). On vérifie aussitôt
que le diagramme (2.5.8.5) est commutatif, et comme $A'$ est un $A$-module
plat, ses lignes sont exactes. Enfin, on a vu

